\documentclass[10pt,twocolumn]{article}
\usepackage[margin=0.72in]{geometry}
\usepackage{booktabs}
\usepackage{graphicx}
\usepackage{hyperref}
\usepackage{xcolor}
\usepackage{enumitem}
\usepackage{microtype}
\newcommand{\sys}{AgentSight Evolution Atlas}
\newcommand{\tool}{\texttt{agent-session-export}}

\title{Process, Outcome, and Survival:\\Visualizing Long-Horizon Agentic Software Evolution}
\author{Anonymous Authors}
\date{}

\begin{document}
\maketitle

\begin{abstract}
Commit-centric software-evolution tools assume that commits are the useful
observation unit and that human contributors retain part of the development
rationale. Long-running coding agents weaken both assumptions: native histories
record second-scale exploration and editing, while no persistent participant
may have observed the full process. Yet native events do not prove durable Git
outcomes, and temporal adjacency does not prove causality or authorship.

We present \sys, an uncertainty-preserving representation and interactive
gallery that keeps three layers distinct: recorded agent process, durable Git
history, and current endpoint state. A Rust exporter normalizes Claude, Codex,
and Gemini histories, follows first-parent Git changes and file lifetimes, and
retains zero, one, or many candidate event--Git associations. A coordinated
TypeScript atlas adapts seven families from thirty years of software-evolution
visualization and adds longitudinal rhythm, process-QA, and resource views.

We evaluate the load-bearing association assumption before using the views for
outcome claims. Exact-hunk evidence passes all controlled held-out gates. On two
mature real-history days, annotators agree on 95.61\% of 933 event--path pairs,
but naturalistic transfer fails prespecified support and calibration gates
(ECE 0.192); endpoint line linkage is also undersupported. We therefore bound
the system to descriptive process, Git, uncertainty, and endpoint views rather
than claiming provenance. The public projection contains 56 sessions, 3,960
path-event rows, 858 path records, 652 Git lifetimes, and 177 commits across six
weeks. In seven local headless-browser repetitions, median first-view visibility
is 234~ms and median two-frame scrub response is 83~ms. These results establish
a working visual instrument and the claim boundary required to use it honestly.
\end{abstract}

\section{Introduction}
Software-evolution visualizations made large histories perceptible by choosing
an observation unit. SeeSoft compressed current lines into pixels
\cite{eick1992seesoft}; Evolution Matrix compared files across versions
\cite{lanza2002evolution}; CodeCity and stable software maps turned repositories
into territories \cite{wettel2008codecity,steinbrueckner2010history,kuhn2010cartography};
code\_swarm and Gource animated commit activity
\cite{ogawa2009codeswarm,caudwell2025gource}; History Flow and Git burndown made
survival visible \cite{viegas2004historyflow,sourced2020hercules}; and forensic
views tied visual structure to maintenance decisions \cite{tornhill2015crime}.
These systems differ visually, but nearly all begin with commits, versions, and
human contributors.

Coding-agent histories expose a different process. A long-running agent may
read a configuration file, search several modules, attempt and revise an edit,
run a test, and later produce a commit. Git retains the durable state transition
but compresses that trajectory. Conversely, a native transcript can record a
path reference without proving that the referenced content entered any commit.
The current tree reveals what remains but discards failed and removed work.
A recent multi-method census likewise finds that commit, pull-request,
configuration, and identity signals recover substantially different agent
populations, making a single repository signal an inadequate ground truth
\cite{khosravani2026census}.

This creates a knowledge-recovery problem. Reviewing ten thousand generated
lines directly does not reveal why the system grew in that shape, while reading
weeks of transcripts is mechanically infeasible. A scalar profiler can account
for semantic units and resources, but structure, temporal rhythm, and
exploration order are not scalars. Visualizations can serve as perceptual units
for these patterns---provided they do not turn weak temporal evidence into
false provenance.

We build \sys{} around that constraint. The system preserves three evidence
layers and coordinates multiple projections with one time, path, session,
vendor, and uncertainty selection state. Rather than propose one novel chart,
we treat the gallery as a visual instrument: established encodings become
comparable lenses over the same event-resolved evidence model. The design
covers line pixels, matrices, stable maps, particle playback, strata, forensic
networks, storylines, and six longitudinal views.

The central scientific risk is the join between native events and actual Git
changes. We therefore evaluate it first and accept a negative transfer result.
Controlled exact-hunk evidence is reliable, but our naturalistic sample does
not satisfy the frozen calibration and null-support thresholds. Consequently,
the deployed atlas exposes candidates and mismatches but does not label an
agent as a Git author, a read as causal, or a current line as agent-derived.

This paper contributes:
\begin{enumerate}[leftmargin=*]
  \item an event-resolved representation that separates recorded process,
  durable Git outcome, and frozen endpoint state while retaining association
  ambiguity and file-lifetime discontinuities;
  \item a coordinated visual atlas covering seven historical software-
  evolution families and longitudinal agent-process views over one stable
  interaction grammar; and
  \item a controlled and double-annotated naturalistic evaluation that both
  validates exact-hunk matching in its supported setting and rules out stronger
  real-history provenance and line-survival claims for the present artifact.
\end{enumerate}

\section{Evidence Model and Requirements}
\subsection{Three Non-Interchangeable Layers}
The \emph{process layer} contains vendor-native sessions, models, timestamps,
tool actions, statuses, repository-relative path references, edit sizes and
body-free fingerprints, and reported token counters. A recorded test or lint
action is evidence that an action occurred, not that every prior edit was
verified or that the result was correct.

The \emph{durable layer} contains first-parent Git commits, parents, pseudonymous
author labels, additions, deletions, and rename-aware changed paths. It is the
authoritative source for integrated repository state transitions. The
\emph{endpoint layer} contains current blobs, file-lifetime survival, and Git
blame origins at one frozen revision. Endpoint survival means only that a Git
file lifetime reaches that revision; it is not a causal or time-to-event claim.

An event--Git association is an explicit edge between layers, never a relabeling
of either endpoint. Eligible write--path pairs retain no candidate, one
candidate, or multiple candidates. Ordered read-before-write edges record
temporal order within a prompt/session, not information flow or causality.
Native vendor labels remain distinct from Git authorship.

\subsection{Lifetime and Censoring Semantics}
A file lifetime begins at Git addition, follows a detected rename, and ends at
deletion. Recreating the same literal path begins a new lifetime. A write during
the deletion gap may target only the subsequent add, preventing a convenient
path name from silently joining two different objects.

Every naturalistic observation requires future Git history. Our frozen primary
retrieval interval is 15 minutes before through 24 hours after an event; the
independent oracle packet extends from 24 hours before the observation day
through seven days after. An observation without that future horizon is
right-censored rather than treated as an unmatched event. This rule excludes
our July~14 cell from every association, calibration, lineage, and transfer
metric.

\subsection{Requirements}
The design follows four requirements. R1 preserves layer identity and missing
evidence. R2 makes uncertainty and disagreement visually inspectable. R3 keeps
repository geography stable during time playback so frames remain comparable.
R4 coordinates every compatible view through shared cursor, interval, vendor,
association-state, path, and session selections. These requirements favor a
two-dimensional multi-view instrument over an immersive 3D city: the latter
adds navigation cost without adding an evidence dimension in this artifact.

\section{System}
\subsection{Exporter and Canonical Artifact}
\tool extends AgentSight's vendor-neutral Rust parser. It accepts an explicit
repository, time interval, optional native session paths, and a frozen Git
revision. Claude, Codex, and Gemini tool schemas are normalized into a common
event IR. Repository-relative paths remain exact; external paths are reduced to
groups. Edit bodies, command bodies, prompts, secrets, native file locations,
and author emails are not serialized. Edit hashes permit controlled hunk
comparison without publishing source text.

Exact-hunk comparison is deliberately limited to native edits with one
repository path and one hunk. Multi-file or multi-hunk patches retain their
path observations but carry no event-level exact-hunk claim.

The Git reader walks first-parent history with 50\% rename detection, records
all changes in the audit window, constructs file lifetimes, and measures
current blob sizes. For an eligible write--path pair, it ranks compatible
changes by exact-hunk evidence, direct versus compatibility path match,
event-relative direction and distance, and commit hash. Prebirth and rename
compatibility are tied. No matching algorithm forces a bijection, so squash,
split, ambiguity, unmatched events, and Git-only changes remain visible.

The canonical \texttt{agentsight.longitudinal.v1} JSON carries complete public
semantics. Normalized event JSONL, Perfetto Trace Event JSON, and Gource custom
logs are also available as explicitly lossy compatibility projections. The
Perfetto projection emits process and Git events on separate tracks; the Gource
projection contains only durable Git changes.

\begin{figure*}[t]
  \centering
  \IfFileExists{../../vis-gallery/artifacts/paper-atlas.png}{%
    \includegraphics[width=\textwidth]{../../vis-gallery/artifacts/paper-atlas.png}%
  }{%
    \fbox{\parbox{0.92\textwidth}{\centering Generate the local atlas plate with \texttt{npm run capture} in \texttt{vis-gallery}.}}%
  }
  \caption{The atlas uses one navigation and playback grammar across every
  family. This field-atlas view overlays recorded read/write rates and Git
  commits while keeping their legends and evidence layers separate.}
  \label{fig:atlas}
\end{figure*}

\subsection{Gallery Projection}
A privacy-checked Python projection derives hourly buckets, deterministic path
coordinates, file-level activity and risk summaries, current-line blame pixels,
survival cohorts, Git co-change edges, ordered read-before-write edges, and Git
ownership aggregates. It rejects serialized prompt, command, preview, old/new
string, content, and absolute home-path fields. The projection preserves the
frozen endpoint revision and marks each source day as mature-descriptive or
right-censored-excluded.

The React/TypeScript client uses ECharts~6.1 for dense canvas charts,
Cytoscape.js~3.34 for networks, uPlot~1.6 for time-series traces, and custom
canvas renderers for SeeSoft pixels and particles. Vite produces a static local
application. A single request-animation-frame clock advances the shared cursor;
filters and selections are typed application state rather than renderer-local
state.

\subsection{Visual Families}
Table~\ref{tab:views} summarizes the visual grammar. Each historical family
contributes a distinct perceptual unit rather than a decorative theme.

\begin{table*}[t]
\centering
\small
\begin{tabular}{p{0.18\textwidth}p{0.30\textwidth}p{0.43\textwidth}}
\toprule
Family & Implemented views & Review question / evidence boundary \\
\midrule
Pixels / SeeSoft & Current-line age pixels; touch and association lanes & Which current lines are old, and which files received recorded attention? Blame age and process activity remain separate. \\
Matrix & File--day evolution; association-state matrix; named signals & Which files burst, pulse, shrink, or remain short-lived, and where is join evidence absent or ambiguous? \\
Maps / cities & Stable treemap; deterministic constellation; directory cartogram & Where did attention move while repository geography remained comparable? Area is current Git size, not agent contribution. \\
Animation & Agent--file particle field; Git commit pulses; recent wake & In what order did recorded activity and durable commits occur? Trails encode recency only. \\
Rivers / strata & Activity river; lifetime cohorts; survival ledger; Git sediment & What activity and file cohorts survive at the frozen endpoint? \\
Crime scene & Hotspot treemap; Git co-change; ordered read-before-write & Where should review start, and which relationships are commit correlation versus temporal process order? \\
Storylines & Session journeys; Git author--directory flow; cast list & How did pseudonymous sessions traverse repository regions, separately from Git ownership? \\
Agent rhythms & Punch card; vital signs; semantic flow; verification lag; receipt; paired days & What clock rhythm, resource distribution, and recorded process-QA lag characterize long runs? \\
\bottomrule
\end{tabular}
\caption{Seven established families plus longitudinal agent-process views.}
\label{tab:views}
\end{table*}

Stable layout is critical to the map family: playback changes brightness and
activity overlays, not coordinates. Matrix rows also remain stable across days.
Forensic views deliberately split Git co-change from ordered process edges.
The storyline view assigns one line to each pseudonymous session but uses a
separate Sankey projection for Git authors. These separations make uncertainty
visible through spatial and color channels rather than explanatory footnotes
alone.

\section{Evaluation}
The completed scientific RQ asks whether native events can be associated with
durable Git and current-line outcomes accurately enough for joined claims. Two
artifact evaluations then report what the atlas exposes descriptively and how
the local implementation responds at its tested scale. The original research
program's behavioral-replication RQ and four-condition human review-utility RQ
remain unexecuted. We do not relabel browser smoke tests or author inspection as
answers to them.

\subsection{RQ1 Protocol}
The controlled corpus contains known target and null pairs plus rename,
delete/recreate, ambiguity, squash, split, moved-and-rewritten line, clock-skew,
merge, and pathless cases for every native schema. Thresholds are calibrated on
60 separate pairs and evaluated once on 100 held-out pairs (50 target and 50
null). A path method must have at least 50 positive and 50 null pairs, Wilson
95\% lower bounds of 0.90 for precision and 0.85 for target recall and null
specificity, and ECE at most 0.10. A line overlay additionally requires 100
linked hunks and a joint-precision lower bound of 0.95.

The naturalistic corpus samples two mature AgentSight observation days. Two
independent annotators see oracle packets constructed independently of method
scores. They label evidence-compatible target, null, or unadjudicable---never
causality or authorship. A third July day remains descriptive because the
endpoint provides only 4.8 future hours. All bootstrap intervals use 10,000
resamples with a frozen seed.

\subsection{RQ1 Results}
On the controlled held-out set, the proposed method classifies all 50 targets
and all 50 nulls correctly. Precision, recall, and specificity are 1.0 with
Wilson lower bounds of 0.929; ECE is 0.030 and Brier score is 0.00135. The
nearest literal-path baseline reaches 0.667 precision and 0.50 null
specificity. Removing exact-hunk evidence reproduces that failure. A
literal-only method with exact hunks matches the proposed method, so the
supported advantage belongs to exact-hunk evidence, not lifetime/rename logic.

The two mature real days contain 933 eligible event--path pairs. Annotators
agree on 95.61\% of labels (Cohen's $\kappa=0.827$); reconciliation yields 110
targets, 45 nulls, and 778 unadjudicable pairs. After call-ID deduplication,
single-file/single-hunk safety, and conservative shell-path parsing, the method
classifies 103/110 targets and 45/45 nulls. Top-1 precision is 0.972, positive
recall is 0.936, classification accuracy is 0.955, the null-specificity Wilson
lower bound is 0.921, and ECE is 0.192. It nevertheless fails the frozen gate:
null support is 49 rather than 50 and ECE exceeds 0.10. The selected mode is
\texttt{descriptive\_only}. Mature days contain Gemini sessions but no eligible
Gemini write--path pairs, so cross-vendor transfer is not established.

For 110 agreed target events, endpoint line linkage emits only 37 predictions:
36 correct and one blame mismatch. Its 0.973 point precision has a 0.862 lower
bound, failing both the 100-link support and 0.95 lower-bound gates. Accordingly,
the SeeSoft view displays real Git-blame ages and process attention, but no
validated event-to-current-line attribution overlay.

RQ1 therefore has a mixed answer. Exact-hunk evidence works in the controlled
protocol; naturalistic calibration, null support, line linkage, and Gemini
transfer are insufficient. This negative boundary prevents stronger but
unsupported ``agent-authored surviving code'' claims.

\subsection{Artifact Experience on Real History}
The public atlas combines 56 deduplicated native sessions, 3,960 path-resolvable
event rows, 858 path records, 652 Git lifetimes, 177 commits, 1,852 Git changes,
and 12,000 current Git-blame line pixels. Of the path records, 693 map to one or
more Git lifetimes; aliases and literal-path reuse make these units non-bijective,
while event-only paths remain distinct. Three observation days span June 2
through July 14, 2026; gaps are retained. The July cell remains process-only.

The same history appears qualitatively different across perceptual units.
The matrix makes path bursts and long quiet intervals visible; the stable map
shows that activity moves without moving repository geography; the activity
river exposes vendor/source volume; lifetime cohorts show survival independently
of attention; and the two forensic graphs make Git co-change and process order
visibly non-equivalent. The named labels---Supernova, Pulsar, White Dwarf,
Dayfly, Nova, Fossil, and Steady---are deterministic heuristics for inspection,
not defect predictors. They supply a vocabulary for hypotheses, not validated
behavioral classes.

This experience supports the representational claim that one canonical model
can drive all seven traditions and expose layer disagreement. It does not show
that reviewers become faster or more accurate. Such a claim requires the
predefined four-condition study (Git-only, event table, joined table, and
coordinated atlas) following visualization-evaluation guidance
\cite{merino2018evaluation}.

\subsection{Local Systems Responsiveness}
We measure the checked 4.82~MB artifact in Playwright Chromium~1.61 at a
1600$\times$1050 viewport. Each of seven repetitions opens a fresh page, visits
all families, and changes the shared range input. Median first-family visibility
is 234~ms (p95 325~ms). Median two-animation-frame scrub response is 83~ms
(p95 133~ms). Median family navigation ranges from 160~ms for strata to 908~ms
for the Cytoscape forensic page; particle animation is 422~ms. The final run
uses 76.6~MB of reported JavaScript heap. Unit tests, a production TypeScript
build, and Playwright tests cover every family, filter state, cursor change, and
representative screenshot.

At this tested scale, the atlas is usable as a local interactive artifact, but
the slower force-layout and particle families identify the next optimization
boundary. These headless measurements establish implementation latency, not
perceptual scalability or human review utility. The 4.77~MB monolithic payload
and 606~KB gzipped JavaScript bundle also motivate lazy family loading and
chunked data access before substantially larger deployments.

\section{Related Work}
\paragraph{Software evolution visualization.}
Our visual vocabulary deliberately descends from SeeSoft
\cite{eick1992seesoft}, Evolution Matrix \cite{lanza2002evolution}, CodeCity and
stable maps \cite{wettel2008codecity,steinbrueckner2010history,kuhn2010cartography},
History Flow \cite{viegas2004historyflow}, animated histories
\cite{ogawa2009codeswarm,caudwell2025gource}, forensic analysis
\cite{tornhill2015crime}, ownership maps \cite{girba2005ownership}, and
storylines \cite{ogawa2010storylines}. Githru provides a strong Git visual-
analytics baseline \cite{kim2021githru}; Hercules operationalizes Git burndown
\cite{sourced2020hercules}. Our novelty lies in a common event-resolved evidence
model and its uncertainty boundary, not in claiming invention of these charts.

\paragraph{Agent process capture.}
RECAP aligns IDE-local Copilot conversation with shadow-repository edits
\cite{he2026recap}; trajectory studies normalize and compare agent action
sequences \cite{majgaonkar2025behaviour,oderinwale2026programs}. Native
cross-vendor histories cover a broader local process but offer weaker evidence
than an instrumented shadow repository, which is precisely why our naturalistic
join is gated rather than assumed.

\paragraph{Survival and lineage.}
Recent work compares agent- and human-associated code survival
\cite{rahman2026survive} and uses AST-aware matching to distinguish deletion
from migration and rewrite \cite{gurov2026clsa}. Our file-lifetime projection
is structural Git evidence. Because our event-to-line stage fails its gate, we
do not infer agent-associated survival from it.

\section{Limitations and Discussion}
The study concerns one repository and local native histories. Missing logs,
clock skew, parallel work, schema differences, rebases, squashes, and private
reasoning constrain observability. Deterministic hashes are pseudonymous
evidence, not strong de-identification. Current Git blame is line-origin
metadata rather than semantic authorship, and first-parent traversal excludes
some branch structure.

The naturalistic sample has 49 adjudicated nulls---one fewer than the frozen
minimum---and no mature Gemini write cell. Its confidence is also miscalibrated.
July's right censoring prevents a convenient but invalid expansion. RQ1 is
therefore a reason to preserve uncertainty, not a validation of the full join.
Named visual patterns remain exploratory, and the browser benchmark does not
demonstrate review benefit.

The deeper value of the atlas is epistemic. Resource profiles make agent work
accountable in semantic units; visual encodings make structural and temporal
patterns perceptible. Long-running agent development needs both: known metrics
for routine control and coordinated views for theory recovery and unexpected
structure. Treating the latter as a scientific instrument requires negative
results and missing evidence to remain visible.

\section{Conclusion}
Long-running coding agents make development more recorded but not automatically
more understandable. \sys{} turns real native histories, Git evolution, and a
frozen endpoint into one coordinated visual atlas while refusing to collapse
their evidence semantics. It demonstrates that thirty years of software-
evolution visual grammar can be reused at event resolution, and it runs
interactively on a real six-week-span artifact. Its evaluation also shows why
visual polish cannot replace measurement: controlled exact-hunk association is
supported, while naturalistic calibration, line linkage, and cross-vendor
transfer are not. The result is a useful descriptive instrument with an
explicit path toward stronger longitudinal studies, not a provenance claim in
disguise.

\bibliographystyle{plain}
\bibliography{references}
\end{document}
